\section{The Two Grants}

The chapter can now state how little it has used. The electron and positron
readings do not require a new stock of primitive objects. They rest on the
two grants already forced by the earlier grammar. The first grant is the
local identity that survived the collapse of truth: the needle, or selection
quotient, by which a reading may stand as the same for a specified purpose.
The second grant is the motion grant: the permission to let the world answer
the stationarity question by fixing which residuals vanish under admissible
variation.

Calling these grants axioms is safe only if the word remains narrow. An
axiom here is not a license to import a hidden universe. It is a bounded
permission without which the grammar cannot continue. The first chapter
showed why unbounded identity was forbidden. The second chapter showed that
finite variation needs a rule by which stationarity can be read. The third
chapter grounded the invariant in a concrete operator. The present chapter
uses exactly those permissions and no third one.

The first grant identifies what would otherwise remain only selected. A
quotient does not say that all differences have vanished. It says that, for
this reading, certain differences no longer count as distinct. The charge
reading needs this grant because a baseline is a selection: it identifies the
reference against which orientation is read. Without the local quotient, the
flat path and the tilted path would be mere names with no stable relation to
the pair residue. With the quotient, each baseline can carry its own
identity while still being distinguished from the other.

The second grant is different. It is not a quotient of presentations but a
world-facing constraint on motion. The grammar may define permitted
variations, actions, boundaries, and residuals, but it still needs the
reading in which stationarity is accepted as the law of the situation. This
grant says that a vanishing residual is not merely a formal convenience; it
is the condition by which the measured path answers. It is the place where
the grammar asks the world for the rule of motion and then carries the
answer as constraint.

The two grants are enough because the charged readings use no additional
privilege. The electron needs the selection of a flat baseline and the
stationary second-variation grammar. The positron needs the selection of a
tilted baseline and the same stationary grammar. Geometry generating gauge
will need boundary identity and motion constraint, not a new substance.
The split will need a carrier discipline that forbids illegal mixing, not a
third metaphysical key. The Cauchy limit and finite forcing will need
extension, compatibility, and residual bounds, all of which belong to the
finite grammar already earned.

This austerity matters because signs are tempting. Once a negative and a
positive reading appear, the reader may want to populate a whole world of
opposed entities by declaration. The grammar forbids that. A sign must be
carried by the invariant, oriented by a baseline, and admitted by a sensor
class. A name that lacks those supports is not yet a reading. The two grants
keep the book from replacing measurement with mythology.

The first grant also prevents the opposite mistake, in which no identity is
allowed and therefore no charge can be read. If the flat baseline could not
be selected as the same baseline across the comparison, the \(-1\) reading
would dissolve. If the tilted baseline could not be selected as its own
comparison, the \(+1\) reading would never become stable. The quotient is
therefore not a philosophical luxury. It is the grammar's way of permitting
orientation to have a reference.

The second grant prevents a merely conventional sign system. If the world
did not answer the stationarity question, a sign could be assigned but not
measured. The charged residue would be a mark in a ledger without a rule
that says which variations are admissible and which residuals must vanish.
The law-of-motion grant bounds the possible readings. It forbids a sign from
being detached from the action that discriminates it.

Together the grants explain why the chapter can be bold without becoming
loose. It reads electron and positron, but it reads them as forced signs of a
single invariant under specified baselines. It lets geometry generate gauge,
but only through boundary grammar. It reaches a continuum reading, but only
through finite conditions. The two grants are the only doors through which
identity enters the chapter. Everything else must be carried, bounded,
quotiented, or forced.

